\caption{\textbf{Sparse versus dense latency tiling, response latency versus anatomical hierarchy, and stimulus locking.} Groups and data as in Fig.~\ref{fig:latency_tiling} (dense tiling, Rastermap clusters 21--28, 4,251 neurons; sparse tiling, clusters 0--7, 44--54 and 61--69, 13,275 neurons).
\textbf{a--c} Noise alone turns dense tiling into sparse-tiling statistics. Dense tiling cells with added noise matched, neuron by neuron, to the sparse-tiling neurons' reliability distribution (noise scaled by $1/\sqrt{\text{trials}}$ per trial type), sorted by their own Rastermap fit on odd trials; \textbf{a} odd, \textbf{b} even trials, six stimulus-aligned types. \textbf{c} Median odd/even reliability, fraction of reliable cells ($r\geq0.3$), cross-validated time--time correlation at lag 0 (odd vs even bins) and peak-latency replication on held-out trial types ($\rho$), for sparse tiling, noisy dense tiling and dense tiling. Error bars, 95\% bootstrap intervals over neurons (1,000 resamples).
\textbf{d} Is sparse tiling's timing less reproducible only because its neurons fire little? Latency replication: Spearman $\rho$ between each neuron's peak latency on the two concordant stimulus-aligned types (odd trials) and on the four other stimulus-aligned types (even trials). Each tiling group is compared with other neurons resampled to have the same firing-rate distribution. At matched rates, sparse tiling replicates less than other neurons (0.18 versus 0.30; 11,570 neurons), whereas dense tiling replicates more (0.75 versus 0.42; 3,949 neurons). Bars, median over 1,000 bootstrap draws; error bars, 95\% interval; $p=0.001$ for both differences (fraction of draws in which the difference reversed).
\textbf{e, f} Response latency does not correlate with the anatomical hierarchy score. Median peak latency per anatomical region against its hierarchy score from tract tracing \cite{harris2019hierarchical}; one dot per anatomical region of cortex (circles) or thalamus (squares) with at least 10 reliable neurons ($r\geq0.3$). Latencies are cross-validated: neurons selected on odd trials, latency read on even trials. \textbf{e} After stimulus onset (55 anatomical regions; Spearman $\rho=-0.25$, permutation $p=0.07$). \textbf{f} In the 150\,ms before movement onset (53 anatomical regions; $\rho=-0.16$, $p=0.24$). Regional latencies are nonetheless reproducible across trial halves (odd versus even trials; $\rho=0.73$ in \textbf{e}, 0.84 in \textbf{f}). \textbf{g} Latency tiling is locked to the stimulus, not to movement. Trials were split into fast and slow reaction times (median 121 and 185\,ms). If responses follow the stimulus, a neuron's fast- and slow-trial responses coincide when aligned to stimulus onset; if they follow movement, they coincide when aligned to movement onset. Bars: correlation between each reliable neuron's fast- and slow-trial responses ($r\geq0.3$; mean over neurons), aligned to stimulus onset (solid; 0--150\,ms after) or to movement onset (hatched; 150\,ms before). Dense tiling 0.52 versus 0.18 (3,265 neurons); sparse tiling 0.16 versus 0.08 (3,082 neurons). Error bars, 95\% bootstrap intervals over neurons; stimulus alignment exceeds movement alignment in all 1,000 resamples for both groups.}
\label{fig:latency_tiling_controls}
